\chapter{\label{chap:chap4}Solução Desenvolvida}

Este capítulo descreve como cada elemento do \textbf{Dashboard Financeiro} foi planejado, construído e validado. A intenção é apresentar a visão integrada da solução final, conectando requisitos funcionais, rastreabilidade técnica e evidências de validação em ambiente real.

\section*{Funcionalidades Implementadas por História de Usuário}

Todas as funcionalidades previstas na fase de planejamento foram entregues e organizadas de acordo com as nove histórias de usuário documentadas no Capítulo~\ref{chap:chap3}. A Tabela~\ref{tab:rastreabilidade-funcionalidades} relaciona cada história (\textbf{US}) à funcionalidade correspondente e aos componentes efetivamente disponibilizados em produção.

\begin{table}[htb]
\centering
\caption{\label{tab:rastreabilidade-funcionalidades}Rastreabilidade entre histórias de usuário e funcionalidades entregues.}
\setlength{\tabcolsep}{0.2cm}
\renewcommand{\arraystretch}{1.25}
\small
\begin{tabular}{|>{\raggedright\arraybackslash}p{3.9cm}|>{\raggedright\arraybackslash}p{4.2cm}|>{\raggedright\arraybackslash}p{6.0cm}|}
\hline
\textbf{História de usuário} & \textbf{Funcionalidade MVP} & \textbf{Implementação entregue} \\
\hline
US-L1 (Larissa — cadastro validado) & Onboarding com validação de CPF e preenchimento automático de endereço & Microserviço ``dashboard-financeiro-cadastro-autenticacao'' com endpoints \texttt{/api/v1/users} e \texttt{/api/v1/auth}, integração ViaCEP, criptografia via \texttt{AuthService} e telas \path{Register.tsx}/\path{Login.tsx}. \\
\hline
US-L2 (Larissa — registros diários) & Cadastro e busca de receitas/despesas categorizadas & Endpoints \texttt{/finance/transactions}, \texttt{TransactionService} com recalculo de saldo, filtros no frontend (\path{Transactions.tsx}) e exportação em tela. \\
\hline
US-L3 (Larissa — metas savings) & Metas do tipo \textit{SAVINGS} com alertas contextuais & \texttt{FinancialGoalController}, cálculo automático de progresso, notificações via Resend e telas \path{Goals.tsx}/\path{GoalAllocationModal.tsx}. \\
\hline
US-R1 (Ricardo — múltiplas contas) & Cadastro e manutenção de contas PF/PJ & Endpoints \texttt{/finance/bank-accounts}, entidade \texttt{BankAccountJpaEntity} e tela \path{BankAccounts.tsx} com formulários modais e paginação. \\
\hline
US-R2 (Ricardo — reconciliação) & Histórico filtrável por categoria, período e conta & Serviços \texttt{TransactionQueryService}, filtros avançados em \path{Transactions.tsx} e componentes de busca rápida no frontend. \\
\hline
US-R3 (Ricardo — limites por categoria) & Metas \textit{EXPENSE\_LIMIT} com monitoramento & Mesmos endpoints de metas, campo \texttt{goalType} diferenciado e alertas configuráveis exibidos em \path{Goals.tsx}. \\
\hline
US-P1 (Patrícia — dashboards analíticos) & Visões Summary/Analytics com tendências e distribuições & \texttt{DashboardService}, componentes \texttt{SummaryView}, \texttt{AnalyticsView} (MUI Charts) e ranking de categorias/metas. \\
\hline
US-P2 (Patrícia — checklists educativos) & Módulo educacional com dicas contextuais & Componentes \path{Education.tsx} e \path{LearningChecklist.tsx} consumindo indicadores do período e exibindo rituais guiados. \\
\hline
US-P3 (Patrícia — notificações externas) & E-mails transacionais sobre metas críticas & Microserviço ``dashboard-financeiro-notificacoes'' (consumo Kafka), templates HTML e registro em \texttt{email\_notification\_log}. \\
\hline
\end{tabular}
\normalsize
\end{table}

Os fluxos foram exercitados pelo roteiro descrito em \texttt{PASSO-A-PASSO.md}, garantindo que cada história fosse validada ponta a ponta (frontend, API Gateway, microserviços e notificações).

\section*{Evidências por História de Usuário}

As subseções a seguir descrevem como cada história de usuário se materializa na interface e nos serviços. Foram reservados espaços para inserir as capturas de tela solicitadas durante a defesa do volume.

\subsection*{US-L1 — Cadastro Validado (Larissa)}
\textbf{Fluxo implementado:} formulário de duas etapas em \texttt{Register.tsx} com máscaras de CPF/CEP, autocomplete ViaCEP e validações de unicidade no endpoint \texttt{/api/v1/users}. O login subsequente utiliza \texttt{/api/v1/auth} e guarda tokens no \textit{local storage}.

\begin{figure}[H]
\centering
\safeincludegraphics[width=0.92\textwidth]{figs/us-l1-step1.png}\\
\footnotesize Formulário inicial (campos vazios).

\par\medskip

\safeincludegraphics[width=0.92\textwidth]{figs/us-l1-step2.png}\\
\footnotesize Dados preenchidos e validação de CPF.

\par\medskip

\safeincludegraphics[width=0.92\textwidth]{figs/us-l1-step3.png}\\
\footnotesize Solicitação do código de verificação.

\par\medskip

\safeincludegraphics[width=0.92\textwidth]{figs/us-l1-email.png}\\
\footnotesize E-mail com o OTP enviado pelo serviço Resend.
\caption{Fluxo US-L1 — onboarding validado (etapas de cadastro e disparo do código).}
\label{fig:us-l1a}
\end{figure}

\begin{figure}[H]
\centering
\safeincludegraphics[width=0.92\textwidth]{figs/us-l1-step4.png}\\
\footnotesize Digitação do código recebido.

\par\medskip

\safeincludegraphics[width=0.92\textwidth]{figs/us-l1-step5.png}\\
\footnotesize Confirmação automática e redirecionamento.

\par\medskip

\safeincludegraphics[width=0.92\textwidth]{figs/us-l1-login.png}\\
\footnotesize Login liberado após validação.

\par\medskip

\safeincludegraphics[width=0.92\textwidth]{figs/us-l1-dashboard.png}\\
\footnotesize Acesso concedido ao painel principal.
\caption{Fluxo US-L1 — confirmação do e-mail e autenticação com redirecionamento ao dashboard.}
\label{fig:us-l1b}
\end{figure}

\subsection*{US-L2 — Registros Diários (Larissa)}
\textbf{Fluxo implementado:} tela \texttt{Transactions.tsx} com formulários modais, categorização colorida, filtros de período e exportação local. O backend recalcula saldos em \texttt{TransactionService} e publica eventos \texttt{transactions.events}.

\begin{figure}[H]
\centering
\safeincludegraphics[width=0.92\textwidth]{figs/us-l2-entry.png}\\
\footnotesize Registro de uma nova receita com categorização e data.

\par\medskip

\safeincludegraphics[width=0.92\textwidth]{figs/us-l2-form.png}\\
\footnotesize Confirmação em tela após o cadastro da operação.

\par\medskip

\safeincludegraphics[width=0.92\textwidth]{figs/us-l2-list.png}\\
\footnotesize Lista filtrada dos últimos 30 dias com a transação lançada.
\caption{Fluxo US-L2 — criação e visualização de transações categorizadas no período de 30 dias.}
\label{fig:us-l2}
\end{figure}

\subsection*{US-L3 — Metas Savings (Larissa)}
\textbf{Fluxo implementado:} cadastro de metas no modo \textit{SAVINGS} em \texttt{Goals.tsx}, cálculo automático de progresso e alertas exibidos no frontend e enviados via Resend quando o status muda.

\begin{figure}[H]
\centering
\safeincludegraphics[width=0.85\textwidth]{figs/us-l3.png}
\caption{Fluxo US-L3 — meta de reserva com barras de progresso e alertas configurados.}
\label{fig:us-l3}
\end{figure}

\subsection*{US-R1 — Contas PF/PJ (Ricardo)}
\textbf{Fluxo implementado:} gerenciamento de múltiplas contas em \texttt{BankAccounts.tsx}, com atributos específicos (apelido, tipo, agência) sincronizados com \texttt{/finance/bank-accounts}.

\begin{figure}[H]
\centering
\safeincludegraphics[width=0.92\textwidth]{figs/us-r1-step1.png}\\
\footnotesize Formulário de contas vazio aguardando o primeiro cadastro.

\par\medskip

\safeincludegraphics[width=0.92\textwidth]{figs/us-r1-step2.png}\\
\footnotesize Conta criada e confirmada com alerta verde.

\par\medskip

\safeincludegraphics[width=0.92\textwidth]{figs/us-r1-step3.png}\\
\footnotesize Edição do registro existente com os botões de ação.
\caption{Fluxo US-R1 — manutenção de contas bancárias PF/PJ, incluindo cadastro e edição de registros.}
\label{fig:us-r1}
\end{figure}

\subsection*{US-R2 — Reconciliação (Ricardo)}
\textbf{Fluxo implementado:} histórico detalhado em \texttt{Transactions.tsx} com filtros combinados (período, categoria, conta) e buscas textuais, apoiado por consultas especializadas no \texttt{TransactionQueryService}.

\begin{figure}[H]
\centering
\safeincludegraphics[width=0.85\textwidth]{figs/us-r2.png}
\caption{Fluxo US-R2 — reconciliação de transações com filtro aplicado aos últimos 30 dias.}
\label{fig:us-r2}
\end{figure}

\subsection*{US-R3 — Limites por Categoria (Ricardo)}
\textbf{Fluxo implementado:} definição de metas \textit{EXPENSE\_LIMIT} em \texttt{Goals.tsx}, painel de acompanhamento e alertas visuais/notificações quando o teto é ultrapassado.

\begin{figure}[H]
\centering
\safeincludegraphics[width=0.85\textwidth]{figs/us-r3.png}
\caption{Fluxo US-R3 — metas \textit{EXPENSE\_LIMIT} com destaque para limites ultrapassados.}
\label{fig:us-r3}
\end{figure}

\subsection*{US-P1 — Dashboards Analíticos (Patrícia)}
\textbf{Fluxo implementado:} visões ``Quadros'' e ``Analytics'' com gráficos de tendência, distribuição por categoria/conta e ranking de metas, alimentadas pelo \texttt{DashboardService}.

\begin{figure}[H]
\centering
\safeincludegraphics[width=0.95\textwidth]{figs/us-p1-1.png}\\
\footnotesize Visão geral com variação do saldo e categorias de maior impacto.

\par\medskip

\safeincludegraphics[width=0.95\textwidth]{figs/us-p1-2.png}\\
\footnotesize Distribuições por categoria e séries temporais de receitas/despesas.

\par\medskip

\safeincludegraphics[width=0.95\textwidth]{figs/us-p1-3.png}\\
\footnotesize Fluxo diário e comparação de saldos por tipo de conta.

\par\medskip

\safeincludegraphics[width=0.95\textwidth]{figs/us-p1-4.png}\\
\footnotesize Consolidação de metas com status e prazos no painel Analytics.
\caption{Fluxo US-P1 — dashboards Summary/Analytics exibindo tendências, distribuições e metas.}
\label{fig:us-p1}
\end{figure}

\subsection*{US-P2 — Checklists Educativos (Patrícia)}
\textbf{Fluxo implementado:} módulo \texttt{Education.tsx} que ativa checklists, mantras e dicas conforme indicadores do período, reforçando os rituais sugeridos durante as mentorias.

\begin{figure}[H]
\centering
\safeincludegraphics[width=0.95\textwidth]{figs/us-p2-1.png}\\
\footnotesize Diagnóstico e foco do momento com saldo educativo do período.

\par\medskip

\safeincludegraphics[width=0.95\textwidth]{figs/us-p2-2.png}\\
\footnotesize Guia rápido conectando módulos (Visão Geral, Transações, Metas, Contas).

\par\medskip

\safeincludegraphics[width=0.95\textwidth]{figs/us-p2-3.png}\\
\footnotesize Laboratório de decisões e mantras de educação financeira.

\par\medskip

\safeincludegraphics[width=0.95\textwidth]{figs/us-p2-4.png}\\
\footnotesize Checklists e rituais aplicados para o ciclo atual.
\caption{Fluxo US-P2 — módulo educacional com diagnósticos, guias e checklists interativos.}
\label{fig:us-p2}
\end{figure}

\subsection*{US-P3 — Notificações Externas (Patrícia)}
\textbf{Fluxo implementado:} eventos de metas e transações geram mensagens no Kafka, consumidas pelo serviço de notificações que envia e-mails via Resend e registra \texttt{email\_notification\_log}.

\begin{figure}[H]
\centering
\safeincludegraphics[width=0.95\textwidth]{figs/us-p3-email.png}
\caption{Fluxo US-P3 — e-mail transacional de boas-vindas orientando primeiros passos e reforçando o uso das notificações.}
\label{fig:us-p3}
\end{figure}

\section*{Arquitetura da Solução}

A arquitetura geral do sistema foi estruturada com base na separação de responsabilidades entre frontend, backend e infraestrutura compartilhada. Para ilustrar o ambiente real em produção, a Figura~\ref{fig:arquitetura-geral} mostra a disposição dos contêineres Docker dentro do Railway, agrupados nos blocos de serviços, infraestrutura e frontend conforme captura da implantação.

\begin{figure}[H]
\centering
\safeincludegraphics[width=0.95\textwidth]{figs/railway.png}
\caption{Ambiente de produção no Railway com os contêineres do \textit{Dashboard Financeiro}.}
\label{fig:arquitetura-geral}
\end{figure}

\section*{Funcionalidades do Sistema}

O sistema \textbf{Dashboard Financeiro} foi projetado para proporcionar ao usuário um ambiente intuitivo e eficiente para o controle de suas finanças pessoais. As funcionalidades foram definidas com base nas entidades modeladas no banco de dados e estruturadas para atender às necessidades práticas do cotidiano financeiro. As principais funcionalidades são:

\begin{itemize}
    \item \textbf{Cadastro de Usuário:} criação de conta com dados pessoais, endereço completo e credenciais de acesso;
    
    \item \textbf{Login Seguro:} autenticação via e-mail e senha, com proteção de rotas por meio do Spring Security~\cite{SPRINGSECURITY2025};
    
    \item \textbf{Gerenciamento de Contas Bancárias:} cadastro de múltiplas contas com tipo (corrente, poupança etc.), número da conta, agência e saldo atual;
    
    \item \textbf{Registro de Transações:} lançamento de entradas e saídas financeiras com valor, data, categoria e tipo (receita ou despesa), atualizando automaticamente o saldo da conta;

    \item \textbf{Classificação por Categoria:} associação de cada transação a uma categoria específica, como alimentação, transporte, moradia, saúde, lazer, entre outras;

    \item \textbf{Visualização de Histórico Financeiro:} exibição do histórico de transações por conta, com filtros por período, tipo e categoria;

    \item \textbf{Metas Financeiras:} definição de objetivos com valor e prazo, permitindo o planejamento e o acompanhamento de metas pessoais;

    \item \textbf{Notificações de Metas:} alertas visuais e e\hyp mails transacionais (via serviço de notificações + Resend) quando metas estão próximas do vencimento ou mudam de status;

    \item \textbf{Painel Resumo:} apresentação gráfica do saldo atual, despesas por categoria, metas ativas e progresso financeiro do usuário;

    \item \textbf{Módulo de Educação Financeira:} orientações, checklists, mantras e rituais diários conectados aos indicadores reais para apoiar novos hábitos;

    \item \textbf{Organização Responsiva:} adaptação da interface a diferentes tamanhos de tela, promovendo usabilidade em computadores e dispositivos móveis;

    \item \textbf{Gerenciamento de Sessão:} controle de login do usuário, com encerramento de sessão e redirecionamento seguro;

    \item \textbf{Segurança dos Dados:} proteção de endpoints, uso de tokens JWT para autenticação e criptografia de senhas no banco de dados~\cite{SPRINGSECURITY2025}.
\end{itemize}

Essas funcionalidades compõem o escopo mínimo viável (MVP) da aplicação, com foco em usabilidade, segurança e utilidade prática. Outras funcionalidades adicionais estão previstas como parte dos trabalhos futuros.

\section*{Descrição da Arquitetura}

A arquitetura full-stack será dividida nas seguintes camadas:

\begin{description}
    \item[Frontend:] Interface desenvolvida em \textbf{React.js}~\cite{REACT2025}, com utilização de bibliotecas como Chart.js ou Recharts para geração de gráficos e Axios para comunicação com a API;
    \item[Backend:] Implementado com \textbf{Spring Boot}, contendo APIs RESTful, controle de autenticação com \textbf{Spring Security}~\cite{SPRINGSECURITY2025}, e persistência dos dados com \textbf{Spring Data JPA};
    \item[Persistência:] Utilização do \textbf{MySQL}~\cite{MYSQL2025}, com configuração via \textbf{Docker} para facilitar testes e ambientes isolados;
    \item[Gerenciamento de Código:] Versionamento com Git e repositório no GitHub.
\end{description}

\section*{Segurança e Privacidade}

A aplicação contará com autenticação baseada em \textbf{JWT (JSON Web Token)}, criptografia das senhas com algoritmos como \texttt{bcrypt} e validações para prevenir injeções e vulnerabilidades comuns no OWASP Top 10~\cite{SPRINGSECURITY2025}. As boas práticas de segurança visam proteger a integridade e a confidencialidade dos dados sensíveis dos usuários.

\section*{Planejamento Metodológico}

O planejamento adotado permaneceu simples e incremental, priorizando o que gerava mais valor em cada ciclo curto. As tarefas foram organizadas diretamente no repositório Git e em anotações compartilhadas com o orientador, com revisões frequentes para validar funcionalidades e ajustar prioridades conforme feedback.

\section*{Metodologia de Desenvolvimento}

A metodologia adotada baseia-se em princípios da engenharia de software moderna e do desenvolvimento ágil. O projeto será dividido em etapas bem definidas:

\begin{itemize}
    \item \textbf{Levantamento de Requisitos:} coleta inicial com base em referências, análise de mercado e objetivos definidos na introdução;
    \item \textbf{Modelagem da Solução:} elaboração de diagramas e definição da arquitetura do sistema;
    \item \textbf{Configuração de Infraestrutura:} preparação do ambiente com Docker, repositórios e integração contínua;
    \item \textbf{Desenvolvimento Incremental:} implementação do sistema por partes (MVP, funcionalidades secundárias, testes);
    \item \textbf{Validação Contínua:} testes de aceitação e com usuários, coleta de feedback e refinamento com base em observações;
    \item \textbf{Revisão e Documentação:} registro das decisões de projeto, escrita do TCC e preparação para apresentação final.
\end{itemize}

Essa metodologia garante alinhamento com os objetivos da solução e permite monitoramento constante da evolução do projeto, com entregas incrementais e foco na qualidade final do produto. Em síntese, o planejamento descrito mantém as entregas técnicas conectadas à proposta de oferecer uma ferramenta educativa, segura e acessível para o controle financeiro pessoal.

\section*{Trilha de Desenvolvimento}

Como o produto está finalizado, esta seção resume a trajetória descrita em \texttt{documentacao-final.md}, detalhando o que foi construído em cada marco:

\begin{itemize}
    \item \textbf{16/09/2025 — Bootstrap do monorepo.} O repositório foi criado com um serviço Spring Boot básico (Gradle + Java 25) e um frontend React/Vite totalmente funcional. Esse commit inicial garantiu scripts de build, hot reload e testes mínimos, permitindo que o time executasse `./gradlew bootRun` e `npm run dev` sem ajustes adicionais.

    \item \textbf{17/09/2025 — Configuração de infraestrutura.} Logo após o bootstrap foram adicionados o \texttt{docker-compose} com SQL Server, Kafka e Zookeeper, além do \texttt{application.yml} parametrizado e da estrutura inicial do Liquibase. Esse passo padronizou ambientes e evitou inconsistências entre máquinas.

    \item \textbf{17/09/2025 — Multiplicação de serviços.} No mesmo dia surgiu o layout completo dos microserviços: domínio financeiro, cadastro/autenticação, notificações e API Gateway. Cada projeto ganhou Gradle Wrapper específico, dependências enxutas, empacotamento individual e documentação de suporte.

    \item \textbf{02/10/2025 — Fluxos de autenticação e primeira tela.} O frontend recebeu contexto de autenticação, formulários de login/cadastro, rotas protegidas e integração com Axios. No backend foram implementados os endpoints JWT (\texttt{/api/v1/auth}, \texttt{/api/v1/users}), filtros de segurança e CORS alinhado ao cliente, abrindo espaço para sessões completas.

    \item \textbf{02/10/2025 — Persistência financeira.} O serviço financeiro ganhou entidades JPA de contas e transações, controllers, serviços, regras de saldo e o \texttt{ApiExceptionHandler}. Esse marco estabeleceu o núcleo transacional e preparou o terreno para dashboards reais.

    \item \textbf{02/10/2025 — Ferramentas de apoio.} Na mesma data foram registrados scripts de migração externa, coleções Insomnia, guias por microserviço e documentação de endpoints, permitindo que novos colaboradores reproduzissem os fluxos rapidamente.

    \item \textbf{20/10/2025 — Camada analítica e eventos.} Entraram categorias, metas, dashboard agregado com \texttt{monthlyTrend}, \texttt{categoryBreakdown} e \texttt{goals}, além de eventos Kafka (\texttt{transactions.events}, \texttt{goals.events}, \texttt{user.events}). Também foi criado o guia operacional \texttt{PASSO-A-PASSO.md}, descrevendo a jornada completa (cadastro → login → contas → metas → dashboard).

    \item \textbf{04/11/2025 — Primeira versão ponta a ponta.} Foram adicionados verificação de e-mail, templates HTML, listener de confirmação, gráficos avançados na tela Analytics, páginas CRUD para contas, categorias, metas, transações e perfil, bem como seeds SQL para demonstração. Essa fase consolidou o MVP completo descrito no Capítulo~\ref{chap:chap3}.

    \item \textbf{10/11/2025 — Preparação para deploy e hardening.} A stack migrou de SQL Server para MySQL, cada serviço recebeu Dockerfile próprio, o compose foi movido para a raiz, implementou-se o client Resend dedicado e foram higienizados segredos, garantindo que a solução pudesse ser provisionada em produção de forma repetível.

    \item \textbf{11/11/2025 — Publicação em nuvem com domínio próprio.} Construí as imagens Docker a partir dos serviços (frontend, gateway, cadastro, financeiro e notificações), subi cada contêiner no Railway, configurei variáveis de ambiente compartilhadas e conectei-os às instâncias gerenciadas de MySQL e Kafka. Em seguida registrei o domínio \path{tcc-leonardopreczevski.com.br} na KingHost, apontei o subdomínio \path{app.tcc-leonardopreczevski.com.br} para o frontend hospedado e validei os registros DNS necessários (SPF, DKIM e tracking) para habilitar envios autenticados via Resend. Nesse marco a aplicação completa ficou disponível publicamente com HTTPS e e-mails de produção.
\end{itemize}

Além desses marcos datados, o período final envolveu ajustes finos no frontend (runtime-config para múltiplos ambientes, componentes de dashboard em MUI Charts, responsividade com Tailwind) e no backend (perfis Docker, scripts de inicialização do MySQL, logs estruturados). A documentação resultante — linha do tempo, guias operacionais e rastreabilidade de funcionalidades — assegura que o leitor compreenda exatamente o que foi construído desde o esqueleto inicial até o produto final hospedável.
